\clearpage
\phantomsection

\addcontentsline{toc}{chapter}{{Mở đầu}}
\chapter*{Mở đầu}

\noindent{\Large \textbf{Bài toán}}

Game hành động Roguelike góc nhìn từ trên xuống đặt ra ba thách thức đồng thời cho tác nhân học tăng cường. Thứ nhất, \textit{phân phối bản đồ thay đổi theo từng episode}: tác nhân không thể ghi nhớ bố cục cố định mà phải khái quát hóa chính sách qua nhiều cấu hình không gian khác nhau. Thứ hai, \textit{không gian hành động đa nhánh $(3,3,3,9)$}: tác nhân phải phối hợp đồng thời hướng di chuyển, quyết định chiến đấu và góc ngắm bắn trong thời gian thực. Thứ ba, \textit{phần thưởng thưa và bị trì hoãn}: các tín hiệu có giá trị cao như tiêu diệt kẻ địch hoặc dọn sạch phòng chỉ xuất hiện sau khi tác nhân hoàn thành chuỗi hành vi liên tiếp gồm tiếp cận, né tránh, căn hướng và tấn công.

Sự kết hợp của ba thách thức này khiến các kỹ thuật đơn lẻ như phần thưởng nội tại (ICM, RND) hay bộ nhớ tuần tự (LSTM) chưa chắc tạo ra cải thiện rõ rệt, từ đó đặt ra câu hỏi: cần kết hợp cơ chế nào để tác nhân học được hành vi chiến đấu hiệu quả trên môi trường sinh theo thủ tục?

\vspace{0.5cm}
\noindent{\Large \textbf{Mục tiêu đề tài}}

Đề tài xác định bốn mục tiêu cụ thể sau, là cơ sở để thiết kế thực nghiệm và đánh giá kết quả xuyên suốt các chương:

\renewcommand{\labelitemi}{$-$}
\begin{itemize}
    \item \textbf{MT1}: Xây dựng tác nhân RL có thể chiến đấu hiệu quả trong môi trường Roguelike thủ tục sinh, đánh giá qua \textit{clear\_rate} và \textit{kills/episode}.
    \item \textbf{MT2}: Giải quyết bài toán phần thưởng thưa và phân phối bản đồ thay đổi bằng cách kết hợp curriculum learning với bản đồ sinh bởi PCGNN.
    \item \textbf{MT3}: Xác định kỹ thuật hỗ trợ nào (ICM, RND, LSTM) thực sự tác động đến hiệu năng thông qua ma trận thí nghiệm loại trừ.
    \item \textbf{MT4}: Kiểm tra giả thuyết về nút thắt biểu diễn quan sát (observation bottleneck) và đánh giá khả năng cải thiện điều hướng của tác nhân mà không cần cài cứng cơ chế tìm đường.
\end{itemize}

\vspace{0.5cm}
\noindent{\Large \textbf{Giải pháp đề xuất}}

Để đạt được các mục tiêu này, đề tài xây dựng một pipeline huấn luyện hoàn chỉnh gồm bốn thành phần chính: (1) PCGNN sinh bản đồ và BFS kiểm tra tính hợp lệ; (2) bộ phân loại độ khó chia bản đồ thành Easy/Medium/Hard và TilemapGenerator nạp vào Unity theo curriculum; (3) CombatAgent với vector quan sát cấu trúc và không gian hành động đa nhánh huấn luyện bằng PPO; (4) mở rộng observation lên 432 chiều bằng local occupancy grid 15$\times$15 trong run62. Ma trận thí nghiệm loại trừ 9 run cho phép cô lập tác động của từng kỹ thuật; run59 và run62 kiểm tra giả thuyết observation bottleneck.

\vspace{0.5cm}
\noindent{\Large \textbf{Câu hỏi nghiên cứu}}

Từ các mục tiêu trên, đề tài xác định hai câu hỏi nghiên cứu chính:

\renewcommand{\labelitemi}{$-$}
\begin{itemize}
    \item \textbf{Q1}: Trong số ICM, RND, LSTM và curriculum learning, kỹ thuật nào tạo ra cải thiện lớn nhất khi huấn luyện agent chiến đấu trong môi trường Roguelike có bản đồ sinh theo thủ tục?
    \item \textbf{Q2}: Khi kết hợp curriculum learning với các kỹ thuật trên, có xuất hiện hiệu ứng cộng hưởng rõ rệt hay không?
\end{itemize}

Để trả lời hai câu hỏi này, đề tài triển khai ma trận thí nghiệm loại trừ gồm chín run chính, trong đó mỗi run thay đổi đúng một biến so với cấu hình tham chiếu. Mỗi cấu hình chỉ được chạy một seed; vì vậy các so sánh được hiểu là xu hướng thực nghiệm, chưa phải kết luận thống kê có khoảng tin cậy.

\vspace{0.3cm}

\noindent{\Large \textbf{Đóng góp của đề tài}}

Thông qua quá trình nghiên cứu, triển khai và thực nghiệm, đề tài đạt được các đóng góp chính sau:

\renewcommand{\labelitemi}{$-$}
\begin{itemize}
    \item Xây dựng pipeline huấn luyện DRL cho game Roguelike trong Unity kết hợp PCGNN, curriculum learning và PPO.
    \item Thiết kế CombatAgent với quan sát 207 chiều và không gian hành động đa nhánh $(3,3,3,9)$, phù hợp cho bài toán di chuyển--ngắm--chiến đấu đồng thời.
    \item Bằng chứng thực nghiệm cho thấy curriculum learning cải thiện phần thưởng trung bình \textbf{+85\%} (2.073 → 3.843) so với PPO baseline, trong khi ICM, RND và LSTM chỉ đạt +16\%--20\%.
    \item Phát hiện entropy plateau 4.40 nat nhất quán trên chín cấu hình ablation, gợi ý mạnh rằng đây là hệ quả của quan sát thiếu đặc trưng không gian chứ không phải giới hạn của PPO hay action space.
    \item Run59 (\emph{path-informed}, +BFS features) cho thấy tìm đường là một nút thắt quan trọng; run62 (432 chiều, 33.45M bước) phá plateau xuống 3.099 nat và đạt clear\_rate 0.626 mà không cần cài cứng BFS.
\end{itemize}

\vspace{0.3cm}

\noindent{\Large \textbf{Bố cục của dự án}}
\vspace{0.5cm}

\renewcommand{\labelitemi}{$-$}
\begin{itemize}
	\item \textbf{Mở đầu}: Bài toán, bốn mục tiêu của đề tài, giải pháp tổng quan và câu hỏi nghiên cứu.
	\item \textbf{Chương 1}: Cơ sở lý thuyết -- RL, PPO, GAE và các kỹ thuật hỗ trợ (ICM, RND, LSTM, PCGNN, curriculum), kèm phân tích hạn chế và cách đề tài xử lý. Giới thiệu đặc trưng của Roguelike làm nền tảng cho thiết kế.
	\item \textbf{Chương 2}: Thiết kế hệ thống -- Tổng quan pipeline end-to-end, sau đó đặc tả MDP, quan sát, không gian hành động, hàm phần thưởng, PCGNN pipeline và curriculum framework.
	\item \textbf{Chương 3}: Cài đặt và thực nghiệm -- Môi trường phần cứng/phần mềm, CombatAgent, cấu hình PPO và ma trận thí nghiệm loại trừ (9 run ablation + run59 + run62).
	\item \textbf{Chương 4}: Kết quả thực nghiệm -- Độ đo và tiêu chí đánh giá, sau đó phân tích kết quả nhóm theo câu hỏi nghiên cứu: kỹ thuật khám phá, curriculum learning, entropy plateau và nút thắt quan sát (observation bottleneck).
	\item \textbf{Kết luận}: Tổng kết kết quả, hạn chế và hướng nghiên cứu tiếp theo.
\end{itemize}
